% GENERATED FILE. Edit canonical lessons, then run npm run generate:books.
% canonical-source-hash: fnv1a64:21326343ea358ebb
% canonical-lessons: UR-C429-nami, UR-C429-thand, UR-C429-jara, UR-C429-sikurna, UR-C429-akarna

\chapter{Damp and Cold}
\label{ch:ur-damp-and-cold}
\hlchaptermodality{429}

\begin{chapteropening}
\textbf{By the end of this chapter:} \emph{I can talk about damp walls, a cold morning and cold hands.}

Complete the last lesson of chapter 429: I can talk about damp walls, a cold morning and cold hands.
\end{chapteropening}

\section[\texorpdfstring{\ur{نمی} (namī)}{نمی (namī)}]{\ur{نمی} (namī) — moisture; damp}

\label{lesson:UR-C429-nami}



\begin{quote}
Before the new one: say the Urdu for steam, then the Urdu for to gather (of clouds).
\end{quote}

\subsection*{You'll want to know: \ur{نمی}}
\textbf{\ur{نمی}} — \emph{namī} — \textquotedblleft{}moisture; damp\textquotedblright{}.

\begin{quote}
\ur{بارش} \ur{کے} \ur{بعد} \ur{دیوار} \ur{میں} \ur{نمی} \ur{ہے۔}
\end{quote}

\begin{quote}
\emph{bārish ke bād dīvār meṅ namī hai.}
\end{quote}

\begin{quote}
After the rain the wall is damp.
\end{quote}

Feminine.

\begin{cousinweb}
\textbf{\ur{نمی}} (\emph{namī}) is built on \emph{nam}, \textquotedblleft{}damp, moist\textquotedblright{}, usually traced to Persian \emph{nam}, with the ending \emph{-ī}.
\end{cousinweb}

\begin{grammarlens}[title={What you've built}]
One more word for getting around.
\end{grammarlens}

\subsection*{Your turn}
\begin{itemize}
\raggedright
  \item \emph{Say it:} \emph{namī}
  \item \emph{Say it:} \emph{namī}, once more
  \item \emph{Say it:} \emph{bārish ke bād dīvār meṅ namī hai.} — After the rain the wall is damp.
  \item \emph{Recall:} say the Urdu for steam, then the Urdu for to gather (of clouds), then say \emph{namī} again
\end{itemize}

\begin{tcolorbox}[breakable,colback=teal!4,colframe=teal!35!black,title={Before you move on}]
What does \ur{نمی} mean? (\textquotedblleft{}Moisture; damp\textquotedblright{}.) Say it once more.
\end{tcolorbox}
\section[\texorpdfstring{\ur{ٹھنڈ} (ṭhanḍ)}{ٹھنڈ (ṭhanḍ)}]{\ur{ٹھنڈ} (ṭhanḍ) — a chill; the cold}

\label{lesson:UR-C429-thand}



\begin{quote}
Before the new one: say the Urdu for to gather (of clouds), then the Urdu for moisture.
\end{quote}

\subsection*{You'll want to know: \ur{ٹھنڈ}}
\textbf{\ur{ٹھنڈ}} — \emph{ṭhanḍ} — \textquotedblleft{}a chill; the cold\textquotedblright{}.

\begin{quote}
\ur{صبح} \ur{کی} \ur{ٹھنڈ} \ur{میں} \ur{ہم} \ur{نے} \ur{چائے} \ur{پی۔}
\end{quote}

\begin{quote}
\emph{subah kī ṭhanḍ meṅ ham ne chāe pī.}
\end{quote}

\begin{quote}
We drank tea in the morning chill.
\end{quote}

Feminine.

\begin{cousinweb}
\textbf{\ur{ٹھنڈ}} (\emph{ṭhanḍ}) is inherited: it belongs with Hindi \emph{ṭhaṇḍā}, \textquotedblleft{}cold\textquotedblright{}, which CDIAL lists under Sanskrit \emph{stabdha}, \textquotedblleft{}stiff, numb\textquotedblright{} (CDIAL 13676).
\end{cousinweb}

\begin{grammarlens}[title={What you've built}]
One more word for getting around.
\end{grammarlens}

\subsection*{Your turn}
\begin{itemize}
\raggedright
  \item \emph{Say it:} \emph{ṭhanḍ}
  \item \emph{Say it:} \emph{ṭhanḍ}, once more
  \item \emph{Say it:} \emph{subah kī ṭhanḍ meṅ ham ne chāe pī.} — We drank tea in the morning chill.
  \item \emph{Recall:} say the Urdu for to gather (of clouds), then the Urdu for moisture, then say \emph{ṭhanḍ} again
\end{itemize}

\begin{tcolorbox}[breakable,colback=teal!4,colframe=teal!35!black,title={Before you move on}]
What does \ur{ٹھنڈ} mean? (\textquotedblleft{}A chill; the cold\textquotedblright{}.) Say it once more.
\end{tcolorbox}
\section[\texorpdfstring{\ur{جاڑا} (jāṛā)}{جاڑا (jāṛā)}]{\ur{جاڑا} (jāṛā) — winter; the cold season}

\label{lesson:UR-C429-jara}



\begin{quote}
Before the new one: say the Urdu for moisture, then the Urdu for a chill.
\end{quote}

\subsection*{You'll want to know: \ur{جاڑا}}
\textbf{\ur{جاڑا}} — \emph{jāṛā} — \textquotedblleft{}winter; the cold season\textquotedblright{}.

\begin{quote}
\ur{جاڑے} \ur{میں} \ur{پہاڑوں} \ur{پر} \ur{برف} \ur{گرتی} \ur{ہے۔}
\end{quote}

\begin{quote}
\emph{jāṛe meṅ pahāṛoṅ par barf girtī hai.}
\end{quote}

\begin{quote}
In the cold season snow falls on the mountains.
\end{quote}

Masculine.

\begin{cousinweb}
\textbf{\ur{جاڑا}} (\emph{jāṛā}) is inherited: CDIAL lists the Hindi cousin \emph{jāṛā}, \textquotedblleft{}coldness, winter\textquotedblright{}, under Sanskrit \emph{jāḍya}, \textquotedblleft{}chilliness\textquotedblright{} (CDIAL 5180).
\end{cousinweb}

\begin{grammarlens}[title={What you've built}]
One more word for getting around.
\end{grammarlens}

\subsection*{Your turn}
\begin{itemize}
\raggedright
  \item \emph{Say it:} \emph{jāṛā}
  \item \emph{Say it:} \emph{jāṛā}, once more
  \item \emph{Say it:} \emph{jāṛe meṅ pahāṛoṅ par barf girtī hai.} — In the cold season snow falls on the mountains.
  \item \emph{Recall:} say the Urdu for moisture, then the Urdu for a chill, then say \emph{jāṛā} again
\end{itemize}

\begin{tcolorbox}[breakable,colback=teal!4,colframe=teal!35!black,title={Before you move on}]
What does \ur{جاڑا} mean? (\textquotedblleft{}Winter; the cold season\textquotedblright{}.) Say it once more.
\end{tcolorbox}
\section[\texorpdfstring{\ur{سکڑنا} (sikuṛnā)}{سکڑنا (sikuṛnā)}]{\ur{سکڑنا} (sikuṛnā) — to shrink; to huddle up}

\label{lesson:UR-C429-sikurna}



\begin{quote}
Before the new one: say the Urdu for a chill, then the Urdu for the cold season.
\end{quote}

\subsection*{You'll want to know: \ur{سکڑنا}}
\textbf{\ur{سکڑنا}} — \emph{sikuṛnā} — \textquotedblleft{}to shrink; to huddle up\textquotedblright{}.

\begin{quote}
\ur{سردی} \ur{میں} \ur{بچے} \ur{کمبل} \ur{میں} \ur{سکڑ} \ur{گئے۔}
\end{quote}

\begin{quote}
\emph{sardī meṅ bachche kambal meṅ sikuṛ gae.}
\end{quote}

\begin{quote}
In the cold the children huddled up under the blanket.
\end{quote}

Past (it …): \ur{وہ} \ur{کرتا} \ur{سکڑ} \ur{گیا}, \emph{voh kurtā sikuṛ gayā}.

\begin{cousinweb}
\textbf{\ur{سکڑنا}} (\emph{sikuṛnā}) is inherited: CDIAL lists the Hindi cousins \emph{sãkaṛnā} and \emph{sikaṛnā}, \textquotedblleft{}to be contracted\textquotedblright{}, under Sanskrit \emph{saṃkuṭati}, \textquotedblleft{}bends oneself together\textquotedblright{} (CDIAL 12825).
\end{cousinweb}

\begin{grammarlens}[title={What you've built}]
One more word for getting around.
\end{grammarlens}

\subsection*{Your turn}
\begin{itemize}
\raggedright
  \item \emph{Say it:} \emph{sikuṛnā}
  \item \emph{Say it:} \emph{sikuṛnā}, once more
  \item \emph{Say it:} \emph{sardī meṅ bachche kambal meṅ sikuṛ gae.} — In the cold the children huddled up under the blanket.
  \item \emph{Recall:} say the Urdu for a chill, then the Urdu for the cold season, then say \emph{sikuṛnā} again
\end{itemize}

\begin{tcolorbox}[breakable,colback=teal!4,colframe=teal!35!black,title={Before you move on}]
What does \ur{سکڑنا} mean? (\textquotedblleft{}To shrink; to huddle up\textquotedblright{}.) Say it once more.
\end{tcolorbox}
\section[\texorpdfstring{\ur{اکڑنا} (akaṛnā)}{اکڑنا (akaṛnā)}]{\ur{اکڑنا} (akaṛnā) — to go stiff; to go numb}

\label{lesson:UR-C429-akarna}



\begin{quote}
Before the new one: say the Urdu for the cold season, then the Urdu for to shrink.
\end{quote}

\subsection*{You'll want to know: \ur{اکڑنا}}
\textbf{\ur{اکڑنا}} — \emph{akaṛnā} — \textquotedblleft{}to go stiff; to go numb\textquotedblright{}.

\begin{quote}
\ur{ٹھنڈ} \ur{سے} \ur{میری} \ur{انگلیاں} \ur{اکڑ} \ur{گئیں۔}
\end{quote}

\begin{quote}
\emph{ṭhanḍ se merī ungliyāṅ akaṛ gaīṅ.}
\end{quote}

\begin{quote}
My fingers went stiff with the cold.
\end{quote}

Past (he …): \ur{وہ} \ur{اکڑ} \ur{گیا}, \emph{voh akaṛ gayā}.

\begin{cousinweb}
\textbf{\ur{اکڑنا}} (\emph{akaṛnā}) is inherited: CDIAL lists the Hindi cousin \emph{akaṛnā}, \textquotedblleft{}to be numbed\textquotedblright{}, under \emph{ākkaḍa}, \textquotedblleft{}cramped, stiff\textquotedblright{}, a form it reconstructs (CDIAL 1013).
\end{cousinweb}

\begin{grammarlens}[title={What you've built}]
One more word for getting around.
\end{grammarlens}

\subsection*{Your turn}
\begin{itemize}
\raggedright
  \item \emph{Say it:} \emph{akaṛnā}
  \item \emph{Say it:} \emph{akaṛnā}, once more
  \item \emph{Say it:} \emph{ṭhanḍ se merī ungliyāṅ akaṛ gaīṅ.} — My fingers went stiff with the cold.
  \item \emph{Recall:} say the Urdu for the cold season, then the Urdu for to shrink, then say \emph{akaṛnā} again
\end{itemize}

\begin{tcolorbox}[breakable,colback=teal!4,colframe=teal!35!black,title={Before you move on}]
What does \ur{اکڑنا} mean? (\textquotedblleft{}To go stiff; to go numb\textquotedblright{}.) Say it once more.
\end{tcolorbox}
